% Generated draft: requires mathematical and source review.
\documentclass[11pt]{article}
\usepackage[margin=25mm]{geometry}
\usepackage{fontspec}
\setmainfont{Latin Modern Roman}
\usepackage{amsmath,amssymb,mathtools,unicode-math}
\setmathfont{Latin Modern Math}
\usepackage{enumitem,needspace,xurl,hyperref}
\setlength{\parindent}{0pt}
\setlength{\parskip}{5pt}
\setlength{\emergencystretch}{3em}
\newcommand{\N}{\mathbb{N}}
\newcommand{\Z}{\mathbb{Z}}
\newcommand{\Q}{\mathbb{Q}}
\newcommand{\R}{\mathbb{R}}
\newcommand{\C}{\mathbb{C}}
\newcommand{\F}{\mathbb{F}}
\newcommand{\A}{\mathbb{A}}
\newcommand{\PP}{\mathbb{P}}
\newcommand{\E}{\mathbb{E}}
\newcommand{\eps}{\varepsilon}
\newcommand{\dd}{\,\mathrm{d}}
\begin{document}
\section*{UnsolvedMath ID: 30004773}
\textbf{Source problem code:} catalogue entry
\subsection{Definitions and mathematical statement}
Let \(\Gamma=(V,E)\) be a finite simple graph, so \(V\) is finite and \(E\subseteq\{\{u,v\}:u,v\in V,\ u\ne v\}\). Let \(p\) range over the odd primes. The supplied record denotes by \(G_\Gamma(\F_p)\) a graph-defined group described as having exponent \(p\) and nilpotency class two. It does not supply a presentation or construction of this group, including the convention relating edges to commutation relations. Consequently, the group notation remains unresolved, and the following target is conditional on recovering its intended definition from the cited source.

For a finite group \(H\), let \(\operatorname{Irr}_{\C}(H)\) denote the set of characters of its irreducible finite-dimensional complex representations; the degree of \(\chi\) is \(\chi(1)\). For each integer \(i\geq1\), define, conditional on the intended construction,
\[
 r_{\Gamma,i}(p)=\#\bigl\{\chi\in\operatorname{Irr}_{\C}(G_\Gamma(\F_p)):\chi(1)=p^i\bigr\}.
\]
The problem is to determine the ordinary representation theory of \(G_\Gamma(\F_p)\) for every finite simple graph \(\Gamma\) and every odd prime \(p\); in particular, for each fixed \(\Gamma\) and every integer \(i\geq1\), determine the function \(p\mapsto r_{\Gamma,i}(p)\) on the odd primes. The supplied question does not assert that this function is polynomial.

The source wording does not clarify whether ``class two'' means class exactly two or at most two, nor how degenerate graphs are accommodated. No restriction excluding such graphs is supplied; no interpretation resolving this issue is imposed here.
\subsection{Short English statement}
For each finite simple graph, determine the complex irreducible representations of its associated group and, in particular, how the number of characters of degree \(p^i\) varies with the odd prime \(p\), for every \(i\geq1\). The supplied record omits the group construction, so a fully specified statement requires source clarification.
\subsection{Sources}
\begin{itemize}
\item Supplied local UnsolvedMath registry record, permanent ID \(30004773\), title: Character Degrees of Graph-Defined Exponent-\(p\) Groups; source problem code \texttt{catalogue entry}. The record belongs to the Oberwolfach Reports Open Problems collection, gives proposed year 2021, and names no proposer.
\item Supplied source citation: \emph{Computational Group Theory (2022)}, Oberwolfach Reports, DOI 10.4171/owr/2021/38, pp. 2027--2087. \url{https://doi.org/10.4171/owr/2021/38}. The literature-source list identifies this as \emph{Oberwolfach Report 38/2021}. Both supplied year designations are retained without reconciliation.
\item Supplied literature reference: Rossmann, \emph{Enumerating conjugacy classes of graphical groups}. \url{https://doi.org/10.1112/blms.12665}.
\item Literature assessment dated \texttt{2026-08-22}: the supplied curation labels the problem open and reports polynomial character-count formulas for special graph families, including edgeless, path, and complete graphs, while describing the general dependence on \(p\) as open. This is supplied dated curation, not an independently verified literature assessment.
\item The registry labels the statement \texttt{corrected\_verified} and reports that the official Oberwolfach report was checked directly and that definitions were restored. These are supplied verification claims; no source visit or independent verification is asserted in this draft, and the supplied record itself contains no construction of \(G_\Gamma(\F_p)\).
\end{itemize}
\end{document}
